%=======================================================================
%  CHAPTER 3 — SEARCH TECHNIQUES  (9 hrs)
%=======================================================================
\section{Search Techniques}

{\color{sub}\footnotesize 9 hrs \gap$\cdot$\gap asked in \mk{all 41 papers}
\gap$\cdot$\gap 3 syllabus sub-topics \gap$\cdot$\gap 392 marks, 11.6\,\% of every mark
ever set \gap$\cdot$\gap worked traces in \emph{Ch3 -- Search Techniques -- Solved
Problems}.}

\lead{Read this first:}
\begin{itemize}
  \item \mk{The chapter is three comparisons and one algorithm.} Learn
        \textbf{BFS Vs DFS} \tS{10}, \textbf{informed Vs uninformed} \tF{4},
        \textbf{A* Vs greedy} \tF{7}, and the \textbf{A* algorithm} itself, and you have
        answered most of what this chapter has ever asked.
  \item \mk{One table carries the whole of \S3.2.} Completeness, optimality,
        time, space, for BFS and DFS. It is asked in \textbf{10 of the 41 papers} and it
        is four rows. Nothing else in the subject pays as well per line learnt.
  \item The other three reliable topics: \textbf{hill climbing's three problems and their
        cures} \tF{6}, \textbf{minimax $+$ alpha-beta} \tF{6}, and
        \textbf{why searching is necessary in AI} \tF{5}, which is a 2-mark opener
        bolted onto a bigger question five times.
  \item \mk{Simulated annealing is asked four times and is not a syllabus
        line.} It only ever appears as \emph{``how does simulated annealing overcome the
        problems of hill climbing''}, so learn it as hill climbing's cure, not as a topic.
  \item Five papers print a \textbf{figure} and ask you to run an algorithm on it: A* on a
        graph (\yr{\textbf{81 Ch}}, \yr{\textbf{80 Ch}}), best-first on a tree
        (\yr{\textbf{\textit{76 Bh}}}), minimax with alpha-beta on a tree
        (\yr{\textbf{79 Ch}}, \yr{76 Ba}). All five are worked in the companion.
        \mk{Two of those papers print the same graph} with different
        heuristics.
  \item Pairs the examiner reliably asks together: \textbf{why search $+$ BFS Vs DFS},
        \textbf{informed Vs uninformed $+$ hill climbing}, and
        \textbf{greedy's drawback $+$ how A* fixes it}.
\end{itemize}

\hr

%-----------------------------------------------------------------------
\T{3.1 Why search, and how a search algorithm is judged}

\Q{Why is searching necessary in AI? On what basis are search algorithms evaluated?}{\m{2}\m{3}\m{4}}
{\color{sub}\footnotesize \tF{5} \yr{72 Ma \m{2}, \textbf{74 Bh} \m{2},
\textbf{78 Ch} \m{2}, 80 Ash \m{2}, 82 Ka \m{3}} \gap$\cdot$\gap
evaluation criteria alone in 3 papers, \yr{79 Ash \m{4}, \textbf{75 Bh} \m{4},
81 Ash \m{2}}, none over 4 marks so no tier chip \gap$\cdot$\gap
\mk{always the 2-mark opener} to a BFS/DFS or hill-climbing question}

\sbs{%
\lead{Why searching is necessary}
\begin{itemize}
  \item Most AI problems have \textbf{no formula} that computes the answer. There is only
        a start state, a set of legal moves and a test for the goal, so the answer has to
        be \textbf{found} rather than calculated.
  \item Once a problem is written as a state space (Ch\,2), \mk{solving it
        \emph{is} searching that space} for a path from initial state to goal.
  \item The space is far too large to enumerate: chess has around $10^{120}$ game paths,
        the 24-puzzle around $10^{25}$ states. \textbf{Brute force is not an option}, so
        the search must be directed.
  \item Search is the mechanism behind almost every other topic on this syllabus: game
        playing (\S3.4), resolution refutation (Ch\,4), rule chaining, planning, and the
        learning algorithms of Ch\,6.
  \item \mk{The one-line answer}: searching is how an agent converts
        knowledge of the rules into a sequence of actions, when it has no direct way to
        compute one.
\end{itemize}}{%
\lead{\mk{The four evaluation criteria}}
\begin{itemize}
  \item \textbf{Completeness} $-$ is the algorithm guaranteed to find a solution when one
        exists?
  \item \textbf{Optimality} $-$ when it finds one, is it guaranteed to be the
        \emph{least-cost} solution?
  \item \textbf{Time complexity} $-$ how many nodes are generated?
  \item \textbf{Space complexity} $-$ how many nodes must be held in memory at once?
  \item \mk{Say what the symbols mean} or the complexities score nothing:
  \begin{itemize}
    \item $b$ $=$ \textbf{branching factor}, the maximum successors of any node;
    \item $d$ $=$ \textbf{depth of the shallowest goal};
    \item $m$ $=$ \textbf{maximum depth} of the search tree, possibly infinite;
    \item $\ell$ $=$ the \textbf{depth limit}, for depth-limited search.
  \end{itemize}
  \item Some papers add \textbf{admissibility} (does it find the optimal path when one
        exists) and \textbf{dominance} (does one heuristic expand fewer nodes than
        another). Mention both in a 4-mark answer.
\end{itemize}}

%-----------------------------------------------------------------------
\T{3.2 Uninformed (blind) search}

\Q{Explain BFS and DFS and compare their performance criteria}{\m{4}\m{7}\m{8}\m{9}}
{\color{sub}\footnotesize \tS{10} \yr{\textbf{72 Ash} \m{9}, 69 Po \m{9},
\textbf{70 Bh} \m{4+4}, 70 Ma \m{1+7}, \textbf{73 Bh} \m{4+5},
\textbf{\texttt{79 Bh}} \m{7}, 79 Jth \m{7}, \textbf{\textit{74 Bh}} \m{8},
\textbf{\textit{71 Bh}} \m{8}, \textit{72 Ma} \m{7}} \gap$\cdot$\gap
\mk{the single most-asked question in the chapter} \gap$\cdot$\gap
\yr{\textbf{\textit{71 Bh}}} and \yr{\textit{72 Ma}} additionally ask
\emph{why} these are called blind}

\sbsr{0.44}{%
\lead{What ``uninformed'' means}
\begin{itemize}
  \item The algorithm is given \textbf{nothing beyond the problem definition}: it can
        generate successors and test for the goal, and that is all.
  \item It has \mk{no estimate of how close a state is to the goal}, so it
        cannot prefer one frontier node over another on merit. Hence \textbf{blind
        search}. That is the answer to ``why are BFS and DFS called blind even though they
        search systematically'': systematic is not the same as informed.
\end{itemize}
\lead{Breadth first search}
\begin{itemize}
  \item Expand \textbf{all nodes at depth $k$ before any at depth $k+1$}.
        \textbf{FIFO queue.}
  \item \checkmark Complete, and \checkmark optimal \emph{when every step costs the same}.
  \item \xmark Memory is the killer: it must hold the whole frontier,
        $O(b^d)$ nodes.
\end{itemize}
\lead{Depth first search}
\begin{itemize}
  \item Expand the \textbf{deepest} unexpanded node; go back only when stuck.
        \textbf{LIFO stack}, or plain recursion.
  \item \checkmark Memory is linear, $O(bm)$: only the current path and its siblings.
  \item \xmark Not complete (it can run down an infinite branch or loop), and
        \xmark not optimal (it returns the first goal found, however deep).
\end{itemize}}{%
\lead{\mk{The comparison table. Learn this.}}
\vspace{-1pt}
\begin{tabularx}{\linewidth}{@{}L{1.75cm}L{2.0cm}L{2.0cm}X@{}}
\toprule
\hrow & \thead{BFS} & \thead{DFS} & \thead{Why} \\
\midrule
Data structure & FIFO queue & LIFO stack & shallowest Vs deepest first \\
Complete & \checkmark yes & \xmark no & DFS can descend forever \\
Optimal & \checkmark yes\newline{\footnotesize (unit step cost)} & \xmark no & BFS meets the shallowest goal first \\
Time & $O(b^d)$ & $O(b^m)$ & $m \geq d$, so DFS can be far worse \\
Space & $O(b^d)$ & $O(bm)$ & \mk{the real difference} \\
Good when & the goal is shallow, memory is plentiful & the tree is deep, memory is tight, solutions are dense & \\
\bottomrule
\end{tabularx}
\vspace{2pt}
\begin{itemize}
  \item \mk{The sentence that earns the last mark}: BFS trades
        \emph{memory} for a guarantee; DFS trades the \emph{guarantee} for memory.
  \item With $b = 10$ and $d = 12$, BFS needs about $10^{12}$ nodes in memory and DFS
        about 120. That is the whole argument in one number.
\end{itemize}}

\penalty0\vspace{2pt}

\sbs{%
\lead{Depth limited search (DLS) \yr{(\textbf{71 Bh} \m{4+4}, 79 Ash \m{4})}}
\begin{itemize}
  \item DFS with a hard cut-off at depth $\ell$: treat any node at depth $\ell$ as having
        no successors.
  \item \textbf{The advantage}, which is what the paper asks: it \mk{fixes
        DFS's infinite-branch failure} while keeping DFS's $O(b\ell)$ memory.
  \item \xmark Still not optimal, and \xmark incomplete in a new way: if
        $\ell < d$ it reports failure even though a solution exists. If $\ell > d$ it is
        not optimal either.
  \item Choosing $\ell$ needs knowledge of the problem. \emph{Eg} for a route between 20
        cities, no path needs more than 19 steps, so $\ell = 19$.
\end{itemize}
\lead{Iterative deepening (IDS) \added}
\begin{itemize}
  \item Run DLS with $\ell = 0, 1, 2, \dots$ until a goal is found.
  \item \mk{Gets BFS's guarantees at DFS's memory cost}: complete, optimal
        for unit steps, $O(b^d)$ time, $O(bd)$ space.
  \item The repeated work is negligible: re-expanding the top of the tree costs a constant
        factor, because the bottom level has more nodes than all the levels above it
        combined.
  \item \textbf{The preferred blind search} when the depth is unknown. Worth one line
        whenever a paper asks for a comparison of blind strategies.
\end{itemize}}

\penalty0\vspace{2pt}

\creamq{Informed Vs uninformed search}{\m{3}\m{4}\m{8}}
{\color{sub}\footnotesize \tF{4} \yr{\textit{74 Ma} \m{8}, \texttt{81 Ba} \m{4+4},
\textbf{\texttt{81 Bh}} \m{3+6}, \textbf{\textit{72 Ash}} \m{8}} \gap$\cdot$\gap
in all four papers it is the \mk{first half of a hill-climbing question}, so
answer it in a table and spend the time on the second half}

\vspace{2pt}
\begin{tabularx}{\textwidth}{@{}L{2.4cm}XX@{}}
\toprule
\hrow & \thead{Uninformed (blind)} & \thead{Informed (heuristic)} \\
\midrule
Extra knowledge & None beyond the problem definition & A \textbf{heuristic} $h(n)$ estimating the cost from $n$ to the goal \\
Chooses by & Position in the queue or stack & The value of an evaluation function \\
Efficiency & Explores blindly, generates many useless nodes & Explores towards the goal, generates far fewer \\
Completeness & BFS, UCS, IDS yes; DFS no & A* yes; greedy and hill climbing no \\
Optimality & BFS and IDS for unit cost, UCS always & A* with an admissible $h$; greedy never \\
Cost & Cheap per node & Needs a heuristic, which costs domain knowledge to design and time to evaluate \\
Examples & BFS, DFS, DLS, IDS, uniform cost & Best first, greedy, A*, hill climbing, simulated annealing \\
\bottomrule
\end{tabularx}
\par\penalty0
\vspace{3pt}

\begin{itemize}
  \item \mk{The line that makes the difference concrete}: on the same
        problem an uninformed search may generate millions of nodes where an informed one
        generates hundreds. The heuristic is not a different algorithm, it is
        \textbf{knowledge added to the same algorithm}.
\end{itemize}

%-----------------------------------------------------------------------
\T{3.3 Informed search: best first, greedy and A*}

\Q{Explain A* search. How does it overcome the drawback of greedy best-first search?}{\m{2}\m{4}\m{6}\m{8}}
{\color{sub}\footnotesize \tF{7} \yr{73 Ma \m{8}, \textbf{\textit{77 Ch}} \m{8},
78 Po \m{8}, \textbf{\textit{73 Bh}} \m{8}, \texttt{82 Ba} \m{2+6}, 78 Ba \m{6},
\textbf{\texttt{82 Bh}} \m{4+5}} \gap$\cdot$\gap
short note on A* alone at \tF{5} \yr{\textit{72 Ma}, 79 Jth, \textbf{\texttt{81 Bh}} \m{4},
\texttt{80 Ba}, 76 Ba \m{3}} \gap$\cdot$\gap
\mk{always phrased as a comparison}, so never explain A* without saying what
greedy gets wrong}

\sbsr{0.5}{%
\lead{The three evaluation functions}
\begin{itemize}
  \item $g(n)$ $=$ the cost \textbf{already paid} to reach $n$ from the start.
  \item $h(n)$ $=$ the \textbf{estimated} cost from $n$ to the goal, the heuristic.
  \item \textbf{Best first} search is the family: keep an \textbf{OPEN} list of generated
        but unexpanded nodes and a \textbf{CLOSED} list of expanded ones, and always
        expand the OPEN node with the best evaluation.
  \begin{itemize}
    \item \textbf{Greedy best first}: $f(n) = h(n)$. Ignores what has already been spent.
    \item \textbf{Uniform cost}: $f(n) = g(n)$. Ignores the goal. This is blind search.
    \item \textbf{A*}: $f(n) = g(n) + h(n)$. \mk{Uses both.}
  \end{itemize}
\end{itemize}
\lead{The A* algorithm, as it should be written}
\begin{itemize}
  \item \textbf{1.} Put the start node on OPEN with $g = 0$, $f = h(\text{start})$.
  \item \textbf{2.} If OPEN is empty, fail.
  \item \textbf{3.} Remove the node of \textbf{lowest $f$} from OPEN, call it $n$, and put
        it on CLOSED.
  \item \textbf{4.} If $n$ is the goal, stop; trace the parent pointers back for the path.
  \item \textbf{5.} Expand $n$. For each successor $m$: compute
        $g(m) = g(n) + \text{cost}(n,m)$ and $f(m) = g(m) + h(m)$; record $n$ as its
        parent; and if $m$ is already on OPEN or CLOSED with a \emph{larger} $g$,
        \mk{replace the old entry and reopen it}.
  \item \textbf{6.} Re-order OPEN by $f$ and go to 2.
\end{itemize}}{%
\lead{\mk{Greedy's drawback, and A*'s fix}}
\begin{itemize}
  \item Greedy asks only \emph{``which node looks nearest to the goal?''}. It therefore
        \textbf{walks into whatever the heuristic flatters}, and a long detour that starts
        with a promising-looking step will be taken in full.
  \item \xmark Greedy is \textbf{not optimal} and \textbf{not complete}: it can loop
        between two nodes that each look good from the other.
  \item A* adds $g(n)$, so \mk{a node is judged on the whole journey, not
        the remainder}. A cheap-looking node that cost a fortune to reach is
        ranked accordingly.
  \item \textbf{The worked case in the companion}: on the \yr{\textbf{80 Ch}} graph,
        greedy takes $S \rightarrow C \rightarrow D \rightarrow F \rightarrow G$ at cost
        \mk{25}, which is the \emph{worst} of the four routes, because
        $h(C) = 5$ looks closest. A* takes $S \rightarrow B \rightarrow E \rightarrow F
        \rightarrow G$ at cost \mk{18}, the optimum.
\end{itemize}
\lead{Comparison in one line each}
\begin{itemize}
  \item \textbf{Greedy}: fast, small memory, no guarantees.
  \item \textbf{A*}: optimal and complete with an admissible $h$, but stores every
        generated node, so \textbf{memory is its limit}.
\end{itemize}}

\penalty0\vspace{2pt}

\creamq{Admissibility and optimality of A*}{\m{4+5}}
{\color{sub}\footnotesize \tP{1} \yr{\textbf{\texttt{82 Bh}} \m{4+5}} \gap$\cdot$\gap
asked as \emph{``define admissibility and optimality; under what conditions is A*
guaranteed to be both''}, so the two \mk{conditions} are the answer}

\sbs{%
\lead{The two properties of a heuristic}
\begin{itemize}
  \item \textbf{Admissible}: $h(n) \leq h^{*}(n)$ for every $n$, where $h^{*}$ is the true
        cost to the goal. The heuristic \mk{never overestimates}; it is
        optimistic.
  \begin{itemize}
    \item Guarantees A* is optimal in \textbf{tree search}.
  \end{itemize}
  \item \textbf{Consistent} (monotone): $h(n) \leq \text{cost}(n,m) + h(m)$ for every edge,
        and $h(\text{goal}) = 0$. A triangle inequality on the heuristic.
  \begin{itemize}
    \item Guarantees A* is optimal in \textbf{graph search}, i.e.\ with a CLOSED list that
          is never revisited. Consistency implies admissibility; the converse is false.
  \end{itemize}
  \item \textbf{Conditions for A* to be admissible and optimal}: $h$ admissible, every
        step cost $> 0$, and the branching factor finite. Add \emph{consistent} if a
        closed list is used.
\end{itemize}}{%
\lead{\mk{Why this is not a pedantic point}}
\begin{itemize}
  \item Both graphs printed in these papers carry heuristics that are
        \textbf{admissible but not consistent}. In \yr{\textbf{81 Ch}},
        $h(B) = 13$ while $\text{cost}(B,E) + h(E) = 6 + 4 = 10$.
  \item The consequence is visible in the trace: A* first reaches $E$ at $g = 12$ through
        $A$, and later finds $E$ at $g = 11$ through $B$. \mk{An
        implementation that refuses to reopen a closed node returns 19 instead of the
        optimal 18.}
  \item So step 5 of the algorithm above, the \emph{replace and reopen} clause, is not
        decoration. Write it, and the trace comes out right.
  \item \textbf{Dominance}: if $h_2(n) \geq h_1(n)$ for all $n$ and both are admissible,
        $h_2$ \emph{dominates} $h_1$ and A* with $h_2$ expands no more nodes. A better
        heuristic is one that is larger while staying under $h^{*}$.
\end{itemize}}

%-----------------------------------------------------------------------
\T{3.4 Hill climbing and simulated annealing}

\Q{Discuss hill climbing, the problems associated with it, and their solutions}{\m{3}\m{4}\m{6}\m{9}}
{\color{sub}\footnotesize \tF{6} \yr{\textbf{69 Bh} \m{9}, 75 Ba \m{6+2},
\texttt{80 Ba} \m{3+5}, \textbf{75 Bh} \m{4}, \textbf{77 Ch} \m{4},
\textbf{\texttt{79 Bh}} \m{4}} \gap$\cdot$\gap
\mk{three problems, three cures}, and the marks are split evenly between
them. Four of these six papers then ask how \textbf{simulated annealing} fixes it}

\sbsr{0.47}{%
\lead{The algorithm}
\begin{itemize}
  \item A \textbf{local search}: keep only the current state, and move to a neighbour that
        improves the objective function. No frontier, no backtracking, no memory of where
        it has been.
  \item \textbf{1.} Evaluate the initial state; if it is the goal, stop.
        \textbf{2.} Generate a neighbour. \textbf{3.} If the neighbour is better, make it
        the current state. \textbf{4.} Repeat until no neighbour is better.
  \item \textbf{Variants}: \emph{simple} hill climbing takes the first better neighbour;
        \emph{steepest ascent} evaluates all neighbours and takes the best;
        \emph{stochastic} picks randomly among the better ones.
  \item \checkmark Almost no memory, and very fast.
        \xmark \mk{Not complete and not optimal}: it finds a
        \emph{local} optimum and stops.
\end{itemize}
\figT{c3_hill_landscape.png}
{\color{sub}\footnotesize The whole answer in one picture: climbing from the current
state reaches the local maximum, not the global one.}}{%
\lead{\mk{The three problems, and their cures}}
\vspace{-1pt}
\begin{tabularx}{\linewidth}{@{}L{1.9cm}XX@{}}
\toprule
\hrow \thead{Problem} & \thead{What happens} & \thead{Cure} \\
\midrule
\textbf{Local maximum} & A peak higher than all its neighbours but lower than the global maximum. Every move looks worse, so the search stops on a suboptimal answer. & \textbf{Backtrack} to an earlier node and try a different direction. Or restart from a new random state (random-restart hill climbing). \\
\textbf{Plateau} & A flat region where all neighbours score the same, so the search has no gradient to follow and wanders at random. & \textbf{Make a big jump} in some direction, to reach a new region of the space. A single step is not enough to get off it. \\
\textbf{Ridge} & A sequence of local maxima along a slope where the ascent runs \emph{across} the ridge rather than along it, so every single-step move goes down. & \textbf{Apply two or more rules at once} before testing, i.e.\ move in several directions at a time. Bidirectional search also helps. \\
\bottomrule
\end{tabularx}
\vspace{2pt}
\begin{itemize}
  \item \mk{Name all three and pair each with its own cure.} A generic
        ``use random restarts'' for all three loses half the marks.
\end{itemize}}

\penalty0\vspace{2pt}

\creamq{Simulated annealing}{\m{4}\m{5}\m{6}}
{\color{sub}\footnotesize \tF{4} \yr{75 Ba \m{2}, \texttt{80 Ba} \m{5},
\texttt{81 Ba} \m{4}, \textbf{\texttt{81 Bh}} \m{6}} \gap$\cdot$\gap
\mk{not a syllabus line}, but asked four times and always as
\emph{``how does it overcome the demerits of hill climbing''}}

\sbs{%
\lead{The idea}
\begin{itemize}
  \item Named after \textbf{annealing} in metallurgy: heat a metal, then cool it slowly so
        its atoms settle into a low-energy crystal instead of freezing into a flawed one.
  \item Hill climbing accepts \emph{only} improving moves, which is exactly why it gets
        stuck. \mk{Simulated annealing sometimes accepts a worse move}, and
        does so less and less as the search proceeds.
  \item A move is picked at random:
  \begin{itemize}
    \item if it \textbf{improves} the state, accept it always;
    \item if it \textbf{worsens} the state by $\Delta E$, accept it with probability
          $P = e^{-\Delta E / T}$.
  \end{itemize}
  \item $T$ is the \textbf{temperature}, lowered on a fixed \textbf{annealing schedule}.
\end{itemize}}{%
\lead{\mk{Why that fixes hill climbing}}
\begin{itemize}
  \item An accepted downhill move is exactly what \textbf{escapes a local maximum}: the
        search can go down one side of the small peak and up the big one.
  \item It also \textbf{crosses a plateau}, since a sideways or slightly worse move costs
        almost nothing in probability.
  \item \textbf{The role of $T$}, which is where the marks are:
  \begin{itemize}
    \item \textbf{High $T$} $\rightarrow$ $P$ near 1, most bad moves accepted, the search
          explores wildly. Nearly a random walk.
    \item \textbf{Low $T$} $\rightarrow$ $P$ near 0, bad moves refused, the search becomes
          \textbf{ordinary hill climbing}.
    \item Starting hot and cooling slowly explores first and refines later.
  \end{itemize}
  \item Note the shape of $e^{-\Delta E/T}$: a \emph{slightly} worse move is much likelier
        to be accepted than a much worse one, at any temperature.
  \item \mk{The theoretical guarantee worth quoting}: if $T$ is lowered
        slowly enough, simulated annealing finds the global optimum with probability
        approaching 1.
\end{itemize}}

%-----------------------------------------------------------------------
\T{3.5 Adversarial search: minimax and alpha-beta}

\Q{Explain the min-max algorithm. What are its drawbacks, and how does alpha-beta pruning solve them?}{\m{2}\m{3}\m{4}\m{6}}
{\color{sub}\footnotesize \tF{6} \yr{71 Ma \m{2+4+3}, \textbf{77 Ch} \m{2+6},
82 Ka \m{6}, 81 Ash \m{6}, \textbf{\textit{76 Bh}} \m{4},
\textbf{\texttt{80 Bh}} \m{2+6}} \gap$\cdot$\gap
short notes on the same at \tF{6} \yr{78 Ba, \textbf{\textit{74 Bh}}, \textit{74 Ma},
\textbf{\textit{72 Ash}}, 78 Po \m{4}, \texttt{82 Ba} \m{3}} \gap$\cdot$\gap
\mk{12 papers between them}, which makes this the second-biggest topic in
the chapter}

\sbs{%
\lead{Minimax}
\begin{itemize}
  \item For a \textbf{two-player, zero-sum, perfect-information} game. One player
        (\textbf{MAX}) maximises the score, the other (\textbf{MIN}) minimises it.
  \item Build the game tree to the terminal states, score each with the utility function,
        then \mk{propagate values upward}:
  \begin{itemize}
    \item at a \textbf{MAX} node, take the \textbf{largest} child value;
    \item at a \textbf{MIN} node, take the \textbf{smallest}.
  \end{itemize}
  \item The root's value is the best MAX can force \textbf{against an opponent playing
        perfectly}, and the move to play is the child that achieves it.
  \item Implemented as a \textbf{depth-first} traversal, so space is $O(bm)$; time is
        $O(b^m)$.
\end{itemize}}{%
\lead{\mk{The drawbacks} (the part the papers ask for)}
\begin{itemize}
  \item \textbf{Exponential time}, $O(b^m)$. Chess has $b \approx 35$ and $m \approx 100$,
        so the complete tree can never be built. \mk{This is the drawback};
        everything else is a consequence.
  \item It \textbf{examines every leaf}, including branches whose outcome cannot possibly
        change the root value.
  \item It assumes a \textbf{perfect opponent}, so it can play too cautiously against a
        weak one.
  \item In practice the tree must be \textbf{cut off at a fixed depth} and the cut-off
        nodes scored by an \textbf{evaluation function}, which makes the answer
        approximate.
\end{itemize}}

\penalty0\vspace{2pt}

\sbs{%
\lead{Alpha-beta pruning}
\begin{itemize}
  \item \mk{Same answer, fewer leaves.} It returns exactly the minimax
        value; it only skips work that cannot affect it.
  \item Carry two bounds down the tree:
  \begin{itemize}
    \item $\alpha$ $=$ the best value \textbf{MAX} can already guarantee. It only rises.
    \item $\beta$ $=$ the best value \textbf{MIN} can already guarantee. It only falls.
  \end{itemize}
  \item \textbf{The two cut rules}:
  \begin{itemize}
    \item at a \textbf{MIN} node, if its value $\leq \alpha$, \textbf{stop} $-$ MAX will
          never choose this branch (\emph{alpha cut-off});
    \item at a \textbf{MAX} node, if its value $\geq \beta$, \textbf{stop} $-$ MIN will
          never allow it (\emph{beta cut-off}).
  \end{itemize}
\end{itemize}}{%
\lead{What pruning buys, and what it does not}
\begin{itemize}
  \item \textbf{Effect}: with a \emph{perfect} move ordering the time complexity falls
        from $O(b^m)$ to $O(b^{m/2})$, i.e.\ \mk{the same time searches twice
        as deep}, or an effective branching factor of $\sqrt{b}$. With random
        ordering it is about $O(b^{3m/4})$.
  \item \xmark It does \textbf{not} change the value returned.
  \item \xmark It does \textbf{not} help in the worst case: with the best move examined
        last everywhere, nothing is pruned and the cost is exactly minimax's.
  \item So it depends entirely on \textbf{move ordering}, which is why real game programs
        invest in ordering heuristics before they invest in searching deeper.
\end{itemize}}

\penalty0\vspace{2pt}

\sbs{%
\lead{How to run alpha-beta in an exam, without losing the thread}
\begin{itemize}
  \item \textbf{Go strictly left to right, depth first.} The examiner's expected answer
        depends on that order.
  \item Write $(\alpha, \beta)$ beside every internal node as you enter it, starting
        $(-\infty, +\infty)$ at the root.
  \item A child \textbf{inherits} its parent's current $\alpha$ and $\beta$.
  \item After each child returns, update the node's own value, then update $\alpha$ (at a
        MAX node) or $\beta$ (at a MIN node).
  \item Test the cut rule \textbf{immediately after each child}, not at the end. Cross out
        the skipped leaves so the pruning is visible.
  \item \mk{State how many leaves were pruned} at the end. It is usually
        worth a mark and it proves the trace was actually run.
\end{itemize}}{%
\lead{Worked in the companion}
\begin{itemize}
  \item \yr{\textbf{79 Ch}}: a 12-leaf, 5-level tree. Minimax value \textbf{3};
        alpha-beta prunes \textbf{2 of the 12 leaves}.
  \item \yr{76 Ba}: a 9-leaf, 2-level tree. Minimax value \textbf{3}; alpha-beta
        prunes \textbf{4 of the 9}, because the first branch returns 3 and sets a high
        $\alpha$ straight away.
  \item \mk{The contrast between the two is the lesson}: the second tree
        prunes twice as large a fraction, purely because its best branch happened to come
        first.
\end{itemize}
\lead{Other adversarial ideas worth a line}
\begin{itemize}
  \item \textbf{Horizon effect}: a fixed-depth cut-off can hide a disaster just beyond the
        last ply, so the agent plays a move that merely postpones it.
  \item \textbf{Quiescence search}: keep searching past the cut-off while the position is
        still volatile.
  \item \textbf{Expectiminimax}: for games with chance, add a \emph{chance} node whose
        value is the expected value over its children.
\end{itemize}}
